\section{Rays and a common renewal arrangement}\label{sec:arrangement}

The radius estimate first reduces a mixed real direction to a mixed
coordinate direction.  Integer heights then let us place two independent
slab sequences in a common system of layers.  We will compare two bounds
for contacts between these sequences.

\subsection{Reduction to a coordinate direction}

\begin{proposition}\label{prop:coordinate-reduction}
Suppose that the annealed walk satisfies
$0<P_0(A_v)<1$ for some $v\in\RR^d\setminus\{0\}$.
There is a coordinate $i$ for which
\[
 P_0(A_{e_i})>0\qquad\text{and}\qquad P_0(A_{-e_i})>0.
\]
\end{proposition}

\begin{proof}
By Lemma~\ref{lem:sign-union}, both signs of $v$ have positive
probability and $P_0(A_v\cup A_{-v})=1$.
Normalize $\|v\|_\infty=1$.
For $\sigma\in\{+,-\}$, let $D_j^\sigma$ be the displacement of
the $j$th slab of the $\sigma v$ sequence, and let $R_j^\sigma$ be
its radius.  Proposition~\ref{prop:radius} and the bound
$|D_j^\sigma|\le R_j^\sigma$ give the mean vector
\[
 b_\sigma=E_\sigma D_1^\sigma,
 \qquad b_\sigma\cdot(\sigma v)=E_\sigma L>0.
\]
In particular $b_\sigma\ne0$.
The strong law gives convergence of the $k$th slab endpoint divided
by $k$ to $b_\sigma$.  Moreover, for each $\varepsilon>0$,
\[
 \sum_{k\ge1}\widehat P_{\sigma v}(R_k^\sigma>\varepsilon k)<\infty,
\]
so $R_k^\sigma/k\to0$ almost surely.  Between the $k$th and
$(k+1)$st endpoints, the displacement from the former is at most
$R_{k+1}^\sigma$.  The number of completed slabs tends to infinity.
Thus the entire trajectory, after its first true cut, has limiting
unit direction
\[
 r_\sigma=\frac{b_\sigma}{|b_\sigma|}.
\]
Lemma~\ref{lem:slabs} transfers this conclusion to the raw law on
$A_{\sigma v}$.  This argument uses slab count rather than elapsed
time and requires no moment of a slab's duration.

The two rays are distinct, since their projections on $v$ have
opposite signs.  The following argument proves the special case of
Simenhaus's antipodality principle needed here~\cite[Proposition~1]{Simenhaus2007}.
We claim that $r_-=-r_+$.  Otherwise
$w=r_++r_-$ satisfies
\[
 w\cdot r_+=w\cdot r_-=1+r_+\cdot r_->0.
\]
The preceding endpoint and radius estimates then imply $A_w$ on
both original directional events, so $P_0(A_w)=1$.
Under $\widehat P_w$, the limiting ray therefore exists almost
surely.  By Lemma~\ref{lem:slabs}, this law is represented by an
iid sequence of finite relative slab words.  Deleting finitely
many words removes a finite time prefix and translates all remaining
positions by a fixed vector.  Since $A_w$ holds, positions tend to
infinity in norm, and these operations preserve the limiting ray.
The ray is consequently a tail variable of the iid word sequence.
The tail zero--one law, applied to a countable determining family
of Borel subsets of the unit sphere, makes it deterministic.
No displacement moment in direction $w$ is needed here.

The raw walk has a first true $w$ cut almost surely, and its
conditional future there has law $\widehat P_w$ by
Lemma~\ref{lem:slabs}.  Its raw limiting ray must therefore be that
same deterministic ray.  This contradicts the positive
probabilities of the two distinct rays $r_+$ and $r_-$.
The claim follows.  Choose a coordinate with nonzero projection
on $r_+$; on $A_v$ and $A_{-v}$ the corresponding coordinate
tends to opposite infinities, proving
Proposition~\ref{prop:coordinate-reduction}.
\end{proof}

It now suffices to rule out a mixed coordinate direction.  Relabel
coordinates so that it is $e_1$, and write $p_\pm=p_{\pm e_1}>0$.
Write $D_\sigma=D_{\sigma e_1}$ and
$\widehat P_\sigma=\widehat P_{\sigma e_1}$ for either sign.
The slab laws $\nu_\pm=\nu_{\pm e_1}$ have positive integer widths;
$E_\pm$ will now denote expectation for these coordinate slab laws.
Lemma~\ref{lem:mean-width} and Proposition~\ref{prop:radius},
applied to this coordinate, give
\begin{equation}\label{eq:coordinate-moments}
 E_\pm L<\infty,\qquad E_\pm R<\infty.
\end{equation}
All heights in the remaining proof are integer coordinate heights.
In particular, a walk started inside a finite coordinate-height
interval exits it almost surely: a path confined to that interval
would have finite height supremum without tending to negative
infinity, contrary to Lemma~\ref{lem:finite-supremum}.
Translation invariance and integration give this exit property
quenched for almost every environment, simultaneously for all
integer barriers and starting sites.  Only a countable intersection
is needed.  This property does not require uniform ellipticity.

\subsection{Exact-cut bridges and renewal tails}

For a sign $\sigma\in\{+,-\}$ define, with starting point zero,
\[
 T_j^\sigma=\inf\{n\ge0:\sigma X_n\cdot e_1=j\},\qquad
 u_j^\sigma=P_0(T_j^\sigma<T_{-1}^\sigma),\qquad
 F_\sigma(j)=\nu_\sigma(L>j),\quad j\ge0.
\]
When only one sign is being discussed we omit $\sigma$.
A relative slab word starts at zero; its raw weight is the function
$W$ defined in Section~\ref{sec:regeneration}.

\begin{lemma}[Exact-cut bridges]\label{lem:exact-bridge}
For either sign, the probability under $\widehat P_\sigma$ of a
slab cut at height $H\ge0$ is $u_H$.  The numbers $u_H$ satisfy
\[
 u_0=1,\qquad 1\ge u_H\ge u_{H+1}\ge p_\sigma.
\]
Conditional on a cut at $H$, the list of slabs through that cut has
law $\mathcal B_H^\sigma$ given by
\begin{equation}\label{eq:exact-bridge-list}
 \mathcal B_H^\sigma(\gamma_1,\ldots,\gamma_m)
 =\frac{1}{u_H^\sigma}\prod_{i=1}^m W(\gamma_i),
 \qquad \sum_{i=1}^m L(\gamma_i)=H.
\end{equation}
Concatenation gives the raw path stopped on its first hit of $H$,
conditional on reaching $H$ before $-1$, in the signed height.
Reversing the order of the list preserves $\mathcal B_H^\sigma$;
every word is still traversed in its original chronological order.
\end{lemma}

\begin{proof}
On $\{T_H<T_{-1}\}$, all departures before $T_H$ have height less
than $H$.  The future non-backtracking event from the fresh endpoint
has probability $p_\sigma$.  Its factor cancels the denominator
in conditioning on $D_\sigma$, giving the cut probability $u_H$.
The hitting events are nested in $H$, and a non-backtracking walk
reaches every height, proving the inequalities.

A finite path first hitting $H$ before $-1$ has a unique slab
decomposition.  Its cuts are the strict record times whose heights
are never undercut before the terminal time, together with the
initial boundary.  Between two successive cuts all intermediate
strict records are invalidated before the second cut.  These are
exactly the qualifying words of Lemma~\ref{lem:slabs}.
Conversely, concatenating any such list of total width $H$ gives
one of these bridges.  Distinct words have disjoint departure-height
intervals, so their raw weights multiply.  Their total mass is
$u_H$, which proves~\eqref{eq:exact-bridge-list} and the bridge
identification.  Reversing the list preserves its total width and
the product of its word weights.  It is a bijection on admissible
lists, proving the last assertion.  At $H=0$ the list is empty.
\end{proof}

\begin{lemma}[Renewal-tail comparison]\label{lem:renewal-tail}
For either sign and all integers $n'\ge n\ge0$,
\begin{equation}\label{eq:renewal-tail}
 0\le u_n-u_{n'}\le(n'-n)F(n).
\end{equation}
Moreover,
\begin{equation}\label{eq:dyadic-width-tail}
 \sum_{l\ge0}2^lF(2^l)<\infty.
\end{equation}
\end{lemma}

\begin{proof}
The slab cut heights form a renewal process with increments $L$.
Partitioning according to its last cut at or before $n$ gives
\[
 \sum_{k=0}^n F(k)u_{n-k}=1.
\]
Subtracting this identity at $n+1$ from the identity at $n$ yields
\[
 \sum_{k=0}^n F(k)(u_{n-k}-u_{n+1-k})=F(n+1).
\]
Every term on the left is nonnegative, and $F(0)=1$.
Thus $u_n-u_{n+1}\le F(n+1)\le F(n)$; telescoping proves
\eqref{eq:renewal-tail}.  Finally
$\sum_{k\ge0}F(k)=EL<\infty$, and monotonicity of $F$ gives
\eqref{eq:dyadic-width-tail} by grouping this sum into dyadic
intervals.
\end{proof}

\subsection{Two tapes placed from a common top}

Berger's backward-path construction also assembles regeneration slabs
from a terminal reference point~\cite[Sections~2--3]{Berger2008}. In the
arrangement below, placement proceeds downwards on both tapes, while
each word retains its chronological orientation.

Take independent iid tapes $(A_i)_{i\ge1}$ and $(B_i)_{i\ge1}$
with respective word laws $\nu_+$ and $\nu_-$.
Their joint law and expectation will be denoted by $\mathbf P$
and $\mathbf E$.  For a relative word $\gamma$ let $d(\gamma)$
be its terminal displacement, and let $a(\gamma)\in\ZZ^{d-1}$
be the last $d-1$ coordinates of that displacement.

Set the positive top endpoint at zero.  Place $A_1$ with its
terminal there, then place $A_2$ with its terminal at the start of
$A_1$, and continue downwards.  Thus the placed copy of $A_i$ is
translated by $-\sum_{r\le i}d(A_r)$.  Place the negative tape
chronologically starting from $-e_1$, so the placed copy of $B_i$
is translated by $-e_1+\sum_{r<i}d(B_r)$.
Only translations are used; the steps inside each word retain
their chronological orientation.

For a site $x$, call $-x_1$ its \emph{progress below the top}, or \emph{depth}.
A word of width $L$ preceded by total tape width $s$ has all its
departures in the progress interval
\begin{equation}\label{eq:departure-progress}
 (s,s+L]\cap\ZZ
\end{equation}
on either tape.  The one-level offset of the negative starting
point is what makes these intervals agree.  A \emph{contact} is
a site that is a departure of both placed tapes.  The last positive
step into the top is axial, so its departure is $-e_1$, the first
negative departure.  There is always a contact at progress one.

Write $H_i^+=\sum_{r\le i}L(A_r)$ and
$H_j^-=\sum_{r\le j}L(B_r)$, with $H_0^+=H_0^-=0$.
Successive positive common values of these two increasing
sequences divide both tapes into \emph{common blocks}.
This is the intersection of two independent renewal processes, a standard
construction discussed, for example, in~\cite{AlexanderBerger2016}.
Lemma~\ref{lem:common-blocks} proves that all such blocks exist and provides the moments
needed for the entropy comparison.

\begin{lemma}[Common blocks]\label{lem:common-blocks}
The paired lists in successive common blocks are iid.  If $K$ is
the total width of one block, then
\begin{equation}\label{eq:common-width-moment}
 \mathbf E K\le\frac{1}{p_+p_-}<\infty.
\end{equation}
Let $V_+$ and $V_-$ be the sums of the actual chronological
lateral displacements of the positive and negative words in that
block.  Then
\begin{equation}\label{eq:common-lateral-moment}
 \mathbf E(|V_+|+|V_-|)<\infty.
\end{equation}
When passing down through this block, the negative-minus-positive
lateral gap increases by
\[
 S=V_++V_-\in\ZZ^{d-1},\qquad \mathbf E|S|<\infty.
\]
\end{lemma}

\begin{proof}
Initially allow the first positive common width $K$ to be infinite.
If it is reached after $i$ positive and $j$ negative words, that
event is determined by those two finite prefixes: their widths
agree, and no earlier positive partial sums agree.  The unused
tails are therefore independent tapes with their original laws,
independent of the completed block.  To see this without any
stopping-time convention, condition separately on each pair
$(i,j)$, factor the law of the unused tails, and sum over $(i,j)$.
Repeating the argument gives the renewal property at every common
cut that exists.

Let $c_n=u_n^+u_n^-$ be the probability that $n$ is a common cut,
including $c_0=1$.  Partitioning according to the last common cut
at or before $n$ gives the elementary identity
\[
 \sum_{i=0}^n c_{n-i}\mathbf P(K>i)=1.
\]
Indeed, after a common cut at $n-i$, the unused pair of tapes has
its original law, and $K>i$ says precisely that there is no further
common cut through $n$.  The identity remains valid if the next
gap is infinite.  Since $c_j\ge p_+p_->0$, it follows that
\[
 (p_+p_-)\sum_{i=0}^n\mathbf P(K>i)\le1.
\]
Letting $n$ tend to infinity proves
\eqref{eq:common-width-moment}.  In particular $K$ is finite
almost surely.  Restarting at each successive common cut now
justifies the infinite iid common-block sequence.

Let $N_+$ be the number of positive words in the first common
block.  The event $\{N_+\ge i\}$ says that none of
$H_1^+,\ldots,H_{i-1}^+$ belongs to the negative renewal set.
It depends only on the preceding positive words and the negative
tape, and is independent of the whole next word $A_i$.
Tonelli's Theorem therefore gives
\[
 \mathbf E\sum_{i=1}^{N_+}R(A_i)
 =\sum_{i\ge1}\mathbf P(N_+\ge i)E_+R
 =\mathbf E N_+\,E_+R
 \le\mathbf E K\,E_+R.
\]
The last inequality uses the integer bound $L\ge1$.
The same argument applies to the negative tape.  Since
$|a(\gamma)|\le R(\gamma)$, this proves
\eqref{eq:common-lateral-moment} and the moment bound for $S$.

Finally, across the block the positive lower anchor is its old
anchor minus $V_+$ in lateral coordinates, whereas the negative
anchor moves by $V_-$.  Their negative-minus-positive difference
increases by $V_++V_-$, as asserted.
\end{proof}

Let $K_j$ be the successive common widths and put
$W_0=0$, $W_j=\sum_{i=1}^jK_i$.  By
\eqref{eq:departure-progress}, both tapes' departures in block
$j$ have progress in $(W_{j-1},W_j]$.  Consequently a contact
always belongs to the same unique common block on both tapes.
Figure~\ref{fig:arrangement} summarizes the two orientations.
\begin{figure}[tbp]
 \centering
 \begin{tikzpicture}[
  x=1cm,y=1cm,>=Stealth,
  every node/.style={font=\small},
  plus/.style={draw=blue!65!black,line width=1pt},
  minus/.style={draw=red!65!black,line width=1pt},
  guide/.style={draw=black!35,densely dashed,line width=.4pt}
]
  \node[anchor=west,font=\small\bfseries] at (0,.85)
    {(a) Common departure bands};
  \node[align=center] at (2.0,.35) {positive\\tape};
  \node[align=center] at (3.9,.35) {negative\\tape};
  \fill[black!4] (.9,-1.3) rectangle (5.1,-2.6);
  \foreach \yy/\lab in {0/0,-1.3/W_1,-2.6/W_2,-3.9/W_3} {
    \draw[guide] (.9,\yy)--(5.1,\yy);
    \node[anchor=east] at (.75,\yy) {$\lab$};
  }
  \foreach \top/\bottom/\j in {0/-1.3/1,-1.3/-2.6/2,-2.6/-3.9/3} {
    \draw[plus,->] (2.0,\bottom+.18)--(2.0,\top-.18);
    \draw[minus,->] (3.9,\top-.18)--(3.9,\bottom+.18);
    \node at (5.45,{(\top+\bottom)/2}) {$K_{\j}$};
  }
  \draw[->,black!70] (-.15,-.30)--(-.15,-3.60);
  \node[rotate=90,anchor=south] at (-.45,-1.95) {progress $-x_1$};
  \node[align=center,font=\footnotesize] at (2.95,-4.35)
    {Arrows indicate orientation, not paths.\\Both tapes are placed downwards.};

  \begin{scope}[xshift=7.0cm]
    \node[anchor=west,font=\small\bfseries] at (-.15,.85)
      {(b) The one-level offset};
    \node at (1.4,.30) {$+$};
    \node at (3.2,.30) {$-$};
    \fill[black!6] (.7,0) rectangle (3.85,-2.7);
    \foreach \yy/\lab in {0/s,-.8/{s+1},-2.7/{s+L},-3.5/{s+L+1}} {
      \draw[guide] (.7,\yy)--(3.85,\yy);
      \node[anchor=east] at (.55,\yy) {$\lab$};
    }
    \draw[plus,->] (1.4,-2.7)--(1.4,0);
    \draw[minus,->] (3.2,-.8)--(3.2,-3.5);
    \fill[blue!65!black] (1.4,-2.7) circle (2.1pt);
    \draw[plus,fill=white] (1.4,0) circle (2.1pt);
    \fill[red!65!black] (3.2,-.8) circle (2.1pt);
    \draw[minus,fill=white] (3.2,-3.5) circle (2.1pt);
    \node[align=center,font=\footnotesize] at (2.28,-1.72)
      {departure\\heights};
    \node[align=center,font=\footnotesize] at (2.1,-4.10)
      {Both words depart only at integer progress\\in $(s,s+L]$; terminal arrivals are excluded.};
  \end{scope}
\end{tikzpicture}

 \caption{The opposed arrangement, schematically and not to scale.
 (a) Common departure bands; arrows indicate chronological orientation,
 while placement proceeds downwards on both tapes. (b) Endpoint progress
 for a word of width $L$ preceded by total width $s$, with the one-level
 offset restored. Filled points are starts and open points are terminal
 arrivals. Word interiors and lateral positions are not depicted;
 the arrows do not assert monotonicity of the words.}
 \label{fig:arrangement}
\end{figure}
For integers $0\le a<b$, denote by $\mathcal C(a,b]$ the event
of a contact in at least one block with index $a<j\le b$.
The quantity to be estimated is
\begin{equation}\label{eq:contact-count}
 m(J)=\sum_{j=1}^J\mathbf P\bigl(\mathcal C(2^j,2^{j+1}]\bigr),
 \qquad J\ge1.
\end{equation}
It is the expected number of dyadic block-index intervals containing
a contact.  Section~\ref{sec:entropy} bounds it above using the
finite first moment of the lateral increment $S$; the remaining
sections obtain an incompatible lower bound from the forced contact
near the top and the renewal-tail estimate.
